% Generated draft: requires mathematical and source review.
\documentclass[11pt]{article}
\usepackage[margin=25mm]{geometry}
\usepackage{fontspec}
\setmainfont{Latin Modern Roman}
\usepackage{amsmath,amssymb,mathtools,unicode-math}
\setmathfont{Latin Modern Math}
\usepackage{enumitem,needspace,xurl,hyperref}
\setlength{\parindent}{0pt}
\setlength{\parskip}{5pt}
\setlength{\emergencystretch}{3em}
\newcommand{\N}{\mathbb{N}}
\newcommand{\Z}{\mathbb{Z}}
\newcommand{\Q}{\mathbb{Q}}
\newcommand{\R}{\mathbb{R}}
\newcommand{\C}{\mathbb{C}}
\newcommand{\F}{\mathbb{F}}
\newcommand{\A}{\mathbb{A}}
\newcommand{\PP}{\mathbb{P}}
\newcommand{\E}{\mathbb{E}}
\newcommand{\eps}{\varepsilon}
\newcommand{\dd}{\,\mathrm{d}}
\begin{document}
\section*{UnsolvedMath ID: 800074}
\textbf{Source problem code:} catalogue entry
\subsection{Definitions and mathematical statement}
A foliation by planes of \(\R^3\) is a foliation \(\mathcal F\) whose leaves are each homeomorphic to \(\R^2\). Its leaf space is the quotient
\[
L(\mathcal F)=\R^3/{\sim_{\mathcal F}},\qquad x\sim_{\mathcal F}y\iff x\text{ and }y\text{ belong to the same leaf},
\]
equipped with the quotient topology. Two such foliations \(\mathcal F\) and \(\mathcal G\) are topologically equivalent if there exists a homeomorphism \(H:\R^3\to\R^3\) carrying every leaf of \(\mathcal F\) onto a leaf of \(\mathcal G\).

\textbf{Requested argument.} Supply an alternative argument for the assertion attributed in the source to Palmeira's theorem that the leaf space determines a foliation by planes of \(\R^3\), avoiding the disputed step asserting that an infinite braid extends to an isotopy of the plane. Under a topological interpretation of ``determines,'' the intended conclusion would have the form
\[
L(\mathcal F)\cong L(\mathcal G)
\quad\Longrightarrow\quad
\mathcal F\text{ and }\mathcal G\text{ are topologically equivalent},
\]
for every pair \(\mathcal F,\mathcal G\) satisfying the hypotheses of the relevant theorem.

\textbf{Unresolved qualifications.} The supplied record does not state those hypotheses, the regularity of the foliations, or the intended category of equivalence. Consequently, the displayed implication records a possible topological meaning of the requested conclusion; it is not an unconditional assertion for all foliations by planes. It asks only for the existence of an equivalence when the leaf spaces are homeomorphic, without requiring a lift of each specified leaf-space homeomorphism. Whether the original theorem requires that stronger conclusion is unspecified. The source reports counterexamples to the infinite-braid extension proposition, but supplies neither its precise formulation nor the counterexamples. That report is retained as a source qualification and is not independently verified here.
\subsection{Short English statement}
Find an alternative argument for the claim attributed to Palmeira that the leaf space determines a foliation of \(\R^3\) by planes, avoiding the infinite-braid extension step disputed by the source. The theorem's precise hypotheses and intended notion of equivalence must be recovered before the target can be stated without qualification.
\subsection{Sources}
\begin{itemize}
\item Local UnsolvedMath registry, supplied complete source record: permanent ID \(800074\), exact source code \texttt{catalogue entry}, title ``Foliations and Diffeomorphism Groups --- Question 15.'' The registry attributes the problem to Nariman and Yazdi, with proposed year 2025, and labels it \texttt{open}; this is a supplied status label, not an independently verified assessment.
\item Nariman and Yazdi, \emph{Problem List on Foliations and Diffeomorphism Groups} (2025), Question 15. Supplied source URL: \url{https://arxiv.org/abs/2503.18274}. The record reports access on 2026-10-16 and extraction by \texttt{arxiv\_latex}; these are supplied provenance details, not a claim of access during this formalization.
\item The question refers to Palmeira's theorem and a proposition in his paper about extending an infinite braid to an isotopy of the plane. No bibliographic identification of that paper, proposition number, or reference for the reported counterexamples is supplied.
\end{itemize}
\end{document}
